\section{Synthesis: why self-improvement does not compound here}
% =====================================================================================
% discussion.tex carries its own \section + \paragraph structure; demote its \section to a
% subsection locally so it nests cleanly under this Synthesis header.
{\let\section\subsection % Discussion: why verified self-training does not compound (distilled from GPT-6-Astra analysis, 2026-09-18).
% Static, hand-checked. \input by figures.tex.
\section{The discovery-to-assimilation bottleneck}

\paragraph{The mechanism: a discovery$\to$assimilation bottleneck (sharpening, not exploration).}
Our two strongest results jointly localise the failure. Best-of-$N$ search beating lineage training at matched
compute (lineage$-$bestof$N=-0.111\,[-0.138,-0.083]$ over 57 independent trajectories), together with the
large pass@$K\!\gg\!$pass@1 gap, shows
that \emph{useful solutions are already present in the sampling tail} but are not reliably produced in one
sample. The limiting pipeline is therefore
\[
\text{rare verified sample}\;\to\;\text{selected SFT example}\;\to\;\text{weight update}\;\to\;\text{reliably reusable gain},
\]
and it is \emph{lossy}: search can exploit a rare discovery directly, whereas training must compress it into a
shared parameterisation, preserve other skills, and have the gain survive later rounds. The lineage-vs-reset
null (inheriting updates does not beat restarting) is direct evidence \emph{against a useful cumulative learning
state}, not merely evidence that individual updates are small. This is the classic self-training signature:
selected-sample training concentrates probability on \emph{already-reachable} successes faster than it expands
what is reachable (cf.\ STaR, ReST, ReST$^{EM}$, expert iteration; and the RLVR capability-boundary debate,
Yue et al.\ 2025).

\paragraph{A formal fixed point --- with its limits stated precisely.}
Under an \emph{idealised} support-preserving rejection-conditioning operator $R$ (unlimited accepted samples,
exact per-task fitting, no cross-task transfer), exact solution-support coverage
$C_0(\pi)=\Pr_x[p_\pi(x)>0]$ is invariant and $R$ is idempotent after one step: $C_0(R\pi)=C_0(\pi)$ and
$R^2\pi=R\pi$. On tasks with positive success probability $R$ places all mass on accepted outputs; on $p{=}0$
tasks no accepted distribution exists. So an idealised verified-SFT loop is a \emph{one-step fixed point in
exact coverage}. \textbf{We are careful about what this does \emph{not} say.} (i) Neural SFT is not this
operator --- shared parameters can generalise and raise the probability of unsampled outputs, so real SFT need
not preserve support. (ii) \emph{Finite-budget} coverage is \emph{not} invariant:
$C_K(\pi)=\mathbb{E}_x[1-(1-p_\pi(x))^K]$ rises sharply under sharpening even when $C_0$ is fixed --- ``pass@$K$
coverage is a fixed point'' is false. (iii) Zero observed successes only bound $p\lesssim 3/K$; our
``unreachable mass'' is more precisely \emph{undiscovered mass at the evaluated budget}. (iv) A training-operator
fixed point is \emph{not} a parameter-space local optimum. We therefore report an \textbf{operational
self-training plateau / protocol-relative fixed point}, never an unqualified ``local optimum.''

\paragraph{Why verified search escapes where training does not.}
Search does not escape the policy's support; it escapes the \emph{one-sample} bottleneck and avoids the
assimilation bottleneck. For success probability $p$, $P(\text{hit in }N)=1-(1-p)^N$ makes extra draws
valuable when $p$ is small, and the verifier keeps the winner. Search (a) uses a discovery directly rather than
compressing it into weights, (b) allocates candidates per-task rather than via one shared update, and (c) does
not permanently contract the proposal policy. The precise statement: \emph{verified search accesses useful
existing tail capability more efficiently than the tested training loop consolidates it}. (This does not settle
lifetime economics --- a training update could amortise over enough future queries; that needs a separate
deployment-horizon study.)

\paragraph{Diversity contraction as a reinforcing lock (associational).}
Selected-sample SFT can drive a loop: current common winners $\to$ accepted set $\to$ higher probability on
those winners $\to$ fewer alternatives in the next accepted set. The loss is \emph{option value} --- rare
decompositions and kernel families that would enable \emph{later} discoveries. If undiscovered useful mass
$U_t$ shrinks, later rounds find fewer novel successes ($1-(1-U_t)^N\le NU_t$) even as pass@1 rises on familiar
ones. Higher retention among compounders and diversity collapse at large scale are \emph{consistent} with this
lock, but we flag them as association, not proof: low diversity can also reflect successful specialisation, and
the cleanest test is a matched-reward, matched-example, diversity-\emph{preserving} intervention.

\paragraph{The inverted-U, mechanistically.} Partition tasks into \emph{undiscovered} ($p{\approx}0$ at budget),
\emph{learnable frontier} ($0<p<0.2$), and \emph{near-ceiling} ($p$ high). Most self-training value comes from
the middle band. Small models: formation failures ($P(A\mid x)=P(F\mid x)P(A\mid F,x)$ with tiny $P(F)$) starve
the accepted set; mid-scale: the frontier is widest, so discoveries become signal (RSI peaks); large scale:
high coverage but little remaining headroom, so drift redistributes mass within already-covered behaviour
(``churn without gain'' --- supported descriptively; ``only superficial late-layer edits'' is \emph{not} yet
established, as parameter distance is not calibrated functional change).

\paragraph{What would break the plateau, and what our data predict.}
Two distinct breakouts are needed: \emph{accumulation} (retain/express more already-discoverable successes ---
could yield lineage$>$reset) and \emph{frontier expansion} (make new strategies operationally reachable ---
required once the frontier is exhausted). A productive loop needs
$\text{new discoveries}+\text{assimilation}+\text{retention}+\text{headroom}$ simultaneously. The single most
informative next experiment is a \textbf{coverage-injection $\times$ diversity-preserving} factorial: inject
stronger-model / repair / search-generated solutions on \emph{undiscovered} tasks and test whether inherited
gains then emerge --- separating ``lack of discoveries'' from ``loss of discoveries after learning.''
Exploration bonuses, novelty-selected accepted sets, prerequisite curricula, off-policy replay, and
search-amplified expert iteration are ranked candidates; each helps only the specific deficit it targets (e.g.\
entropy cannot fix formation; off-policy replay of redundant winners adds no frontier information).

\paragraph{Scope --- claims we do \emph{not} make.} We do not claim the models' reachable set is mathematically
invariant, that the weights sit at a genuine local optimum, or that RL/RLVR cannot expand capability. We claim,
with \emph{moderate} evidence for the null and a significant positive respectively: under verified
rejection-sampling SFT at the tested scales/budget, carrying the lineage forward does not beat a matched reset
(trajectory-level TOST equivalence at $\delta{=}0.05$, $\mathrm{BF}_{01}{=}5.6$ over 57 independent
trajectories); verified search is a better use of the same compute (CI excludes zero); and the binding internal signals are coverage/formation (small), a mid-scale learnable
frontier, and drift-without-gain plus diversity contraction (large). The synthesis: \emph{expert iteration
without sufficient expert amplification, plus incomplete retention of the apprentice's useful tail.}
}

% =====================================================================================
\section{The accumulated intuition (what three weeks taught us)}
% =====================================================================================
This section is the part that does not fit in a conference paper: the qualitative model we now carry in our
heads about RSI in this domain.

\begin{enumerate}[leftmargin=1.6em]
\item \textbf{Self-training is a sharpening operator, not an exploration operator.} It concentrates probability
on already-reachable successes faster than it expands what is reachable. Everything else follows from this.

\item \textbf{Coverage is the currency, and it is not self-generated.} A round of self-training can only reuse
solutions the model already finds; it cannot manufacture coverage on tasks where the per-sample success
probability is effectively zero at the evaluated budget. Compounding is coverage-limited, and coverage is
exogenous to the loop.

\item \textbf{The bottleneck is discovery$\to$assimilation, not knowledge.} The model represents correctness it
cannot reliably generate (probe result), and search --- which uses discoveries directly --- beats training,
which must assimilate them into weights and retain them. The lossy step is assimilation-and-retention.

\item \textbf{``Operational plateau,'' not ``local optimum.''} We are careful: the null is a protocol-relative
fixed point of the tested verified-SFT loop, not a claim that the weights sit at a parameter-space local
optimum. Finite-budget coverage is \emph{not} invariant; a better loop could still move.

\item \textbf{Diversity contraction is a plausible reinforcing lock.} Selected-sample SFT can raise probability
on current winners at the cost of the rare decompositions that would enable \emph{later} discoveries. We flag
this as association, not proof --- the clean test is a diversity-preserving intervention.

\item \textbf{The open contribution is a better \emph{harness}, not a better model.} Because RSI is a protocol
we run, the way to break the plateau is to change the protocol: inject coverage on undiscovered tasks, preserve
diversity in the accepted set, amplify the ``expert'' with search before assimilation. This is exactly what the
public benchmark's harness-submission track exists to surface.
\end{enumerate}

% =====================================================================================
\section{What would break the plateau: predictions}
% =====================================================================================
Two distinct breakouts are needed and they target different deficits: \emph{accumulation} (retain and express
more already-discoverable successes --- could make lineage finally beat reset) and \emph{frontier expansion}
(make genuinely new strategies operationally reachable --- required once the current frontier is exhausted). A
productive loop needs new discoveries $+$ assimilation $+$ retention $+$ headroom \emph{simultaneously}. Our
single most informative next experiment was, at the time of writing, a \textbf{coverage-injection $\times$
diversity-preserving} factorial. The instrument work in Section~\ref{sec:instrument} has since superseded that.
The coverage-injection control was run, and its injected-task count falls to zero by round~2 because the student
solves the entire held set; injection cannot separate ``lack of discoveries'' from ``loss of discoveries'' on a
task the model has already saturated.

The informative experiment is now the one a repaired grader makes possible for the first time: run the
compounding protocol on \textbf{kernel authoring} rather than correctness-preserving rewriting. Under the
published prompt, models attempt a custom kernel in \textbf{0 of 16} samples; when asked explicitly, \textbf{14
of 16}, of which \textbf{1} verifies. A success rate in that band is far from both the floor and the ceiling
that defeated every metric we tried, which is precisely the regime a compounding study requires. Whether
self-training compounds on a task the model has \emph{not} saturated is, as far as we can tell, still open ---
and it is a question this benchmark could not previously have asked.

% =====================================================================================
\section{Limitations and scope}
% =====================================================================================
We claim, with moderate evidence ($\mathrm{BF}_{01}=5.6$ at the trajectory level): under verified rejection-sampling SFT at the tested scales and budgets, carrying
the lineage forward does not beat a matched reset; verified search is a better use of the same compute; and the
binding internal signals are coverage/formation (small scale), a mid-scale learnable frontier, and
drift-without-gain plus diversity contraction (large scale). We do \emph{not} claim the reachable set is
mathematically invariant, that the weights sit at a genuine local optimum, or that RL/RLVR (as opposed to
rejection-sampling SFT) cannot expand capability. All primary intervals are now computed at the
\emph{trajectory} level rather than the round level, since rounds within a run share a base checkpoint, a seed,
a held-out split and an accumulated adapter; this is the honest unit of replication and it costs us evidential
strength, moving the equivalence Bayes factor from $\mathrm{BF}_{01}=14.7$ (pooling 228 rounds) to $5.6$
(57 trajectories) --- strong to \emph{moderate}. The verdicts survive the correction, and the
search-beats-training positive strengthens under it, because that contrast is made within a round. Per-trajectory
(2--3 seed) uncertainty remains the binding constraint. Finally, all results here are A100; our H100 runs are
reported separately and are non-comparable on two axes at once: different roofline constants, and a harness
that silently rejected every non-PyTorch submission (Section~\ref{sec:instrument}). An earlier version of this
sentence blamed a stronger \texttt{torch.compile} baseline; that attribution was wrong --- the default scorer
never used the compiled baseline --- and is retracted in Section~\ref{sec:instrument}.

% =====================================================================================
\section{Conclusion}
% =====================================================================================
Three weeks of clean-verifier, compute-matched experiments across five rungs of recursion and scales from 0.5B
to 14B converge on one picture: in GPU-kernel optimization, verified self-training \emph{sharpens what a model
already covers but does not explore}, so improvement does not compound --- not because the models lack
capability, but because the loop's feedback channel assimilates rare discoveries lossily and cannot manufacture
the coverage it would need to keep going. The result is a strong, honest null flanked by two positive findings
(search beats training; the correctness wall is a sampling artifact) that together localize the failure and
point at the intervention most likely to break it. KernelAscent ships that intervention as an open, automated
submission track, so the community can try to turn the null into a curve.
